\expandafter\gdef\csname SNside1\endcsname{\node[anchor=north west,inner sep=0pt,text width=48mm,font=\SideFont,align=left] (SNsidebox) at ([xshift=138mm,yshift=-402.91pt]current page.north west) {\textcolor{SNink}{\textbf{课程连接 · 系统边界先于分层}}\par\vspace{1.5mm}操作系统研究一台主机内部如何管理资源；计网研究不同自治主机如何协作。后续遇到``缓存满''时，先判断缓存属于网卡、路由器还是接收进程，才能选对控制范围。};\pgfextracty{\SNsideheight}{\pgfpointdiff{\pgfpointanchor{SNsidebox}{south}}{\pgfpointanchor{SNsidebox}{north}}}\typeout{SNHEIGHT:1:0:\the\SNsideheight}}
\expandafter\gdef\csname SNside2\endcsname{\node[anchor=north west,inner sep=0pt,text width=48mm,font=\SideFont,align=left] (SNsidebox) at ([xshift=138mm,yshift=-95.00pt]current page.north west) {\textcolor{SNink}{\textbf{最小反例 · 星形也可能共享}}\par\vspace{1.5mm}三台主机各用一根线连接集线器，物理图是星形；内部仍向其他端口重复信号，三站共享冲突域。把中心换成交换机后并行能力改变，外形本身没有证明介质独占。\par\smallskip\begin{tikzpicture}[x=1mm,y=1mm,>=Stealth,node distance=5mm]
\node[draw=SNink!50,fill=SNblue,rounded corners=1mm,text width=34mm,align=center,inner sep=1.2mm,font=\fontsize{7.5}{9}\selectfont] (m0) at (0,0) {A};
\node[draw=SNink!50,fill=SNblue,rounded corners=1mm,text width=34mm,align=center,inner sep=1.2mm,font=\fontsize{7.5}{9}\selectfont] (m1) at (0,-11) {集线器};
\node[draw=SNink!50,fill=SNblue,rounded corners=1mm,text width=34mm,align=center,inner sep=1.2mm,font=\fontsize{7.5}{9}\selectfont] (m2) at (0,-22) {B};
\node[draw=SNink!50,fill=SNblue,rounded corners=1mm,text width=34mm,align=center,inner sep=1.2mm,font=\fontsize{7.5}{9}\selectfont] (m3) at (0,-33) {C};
\draw[->,SNink] (m0) -- (m1);
\draw[->,SNink] (m2) -- (m1);
\draw[->,SNink] (m3) edge[bend left=40] (m1);
\end{tikzpicture}
\par{\fontsize{7.5}{9}\selectfont 星形接线，四站仍共享介质}};\pgfextracty{\SNsideheight}{\pgfpointdiff{\pgfpointanchor{SNsidebox}{south}}{\pgfpointanchor{SNsidebox}{north}}}\typeout{SNHEIGHT:2:1:\the\SNsideheight}}
\expandafter\gdef\csname SNside3\endcsname{\node[anchor=north west,inner sep=0pt,text width=48mm,font=\SideFont,align=left] (SNsidebox) at ([xshift=138mm,yshift=-95.00pt]current page.north west) {\textcolor{SNink}{\textbf{课程连接 · 像函数调用的接口}}\par\vspace{1.5mm}同机上层请求下层服务可类比函数调用：约定接口，不需要知道实现细节。两端对等层的协议则约束交换消息的格式与动作，不能把本机调用关系画成远端数据直接穿越介质。};\pgfextracty{\SNsideheight}{\pgfpointdiff{\pgfpointanchor{SNsidebox}{south}}{\pgfpointanchor{SNsidebox}{north}}}\typeout{SNHEIGHT:3:4:\the\SNsideheight}}
\expandafter\gdef\csname SNside4\endcsname{\node[anchor=north west,inner sep=0pt,text width=48mm,font=\SideFont,align=left] (SNsidebox) at ([xshift=138mm,yshift=-95.00pt]current page.north west) {\textcolor{SNink}{\textbf{错项排除 · 确认不自动等于可靠}}\par\vspace{1.5mm}无线链路ACK只确认当前一跳，后续路由仍可能丢包；CRC只检出部分差错。题目出现``确认''``检错''时，先核对其作用端点，再排除把一项局部机制扩大为端到端保证的说法。};\pgfextracty{\SNsideheight}{\pgfpointdiff{\pgfpointanchor{SNsidebox}{south}}{\pgfpointanchor{SNsidebox}{north}}}\typeout{SNHEIGHT:4:5:\the\SNsideheight}}
\expandafter\gdef\csname SNside5\endcsname{\node[anchor=north west,inner sep=0pt,text width=48mm,font=\SideFont,align=left] (SNsidebox) at ([xshift=138mm,yshift=-238.60pt]current page.north west) {\textcolor{SNink}{\textbf{不变量 · 模型与实例分开}}\par\vspace{1.5mm}OSI列出的某层可用功能，不等于该层每一种实际协议都实现了全部功能。IP使用无连接数据报、TCP提供可靠字节流、UDP不提供可靠交付，必须先固定题目给的是哪种模型。};\pgfextracty{\SNsideheight}{\pgfpointdiff{\pgfpointanchor{SNsidebox}{south}}{\pgfpointanchor{SNsidebox}{north}}}\typeout{SNHEIGHT:5:6:\the\SNsideheight}}
\expandafter\gdef\csname SNside6\endcsname{\node[anchor=north west,inner sep=0pt,text width=48mm,font=\SideFont,align=left] (SNsidebox) at ([xshift=138mm,yshift=-95.00pt]current page.north west) {\textcolor{SNink}{\textbf{课程连接 · 流水与CPU同一个计数}}\par\vspace{1.5mm}分组流水和指令流水都应分成``填满流水''与``稳定间隔''。等速链路中首组付出多跳发送代价，之后按瓶颈服务间隔完成；直接把首组总延时乘组数，会把已经重叠的时间重复计算。};\pgfextracty{\SNsideheight}{\pgfpointdiff{\pgfpointanchor{SNsidebox}{south}}{\pgfpointanchor{SNsidebox}{north}}}\typeout{SNHEIGHT:6:7:\the\SNsideheight}}
\expandafter\gdef\csname SNside7\endcsname{\node[anchor=north west,inner sep=0pt,text width=48mm,font=\SideFont,align=left] (SNsidebox) at ([xshift=138mm,yshift=-95.00pt]current page.north west) {\textcolor{SNink}{\textbf{最小反例 · 有连接仍可共享}}\par\vspace{1.5mm}两条虚电路沿同一物理链路依次发送分组，各节点维护不同标识映射即可区分。它们都已建立逻辑连接，却无需各占一根线，这就排除了``虚电路必预留固定带宽''的绝对说法。};\pgfextracty{\SNsideheight}{\pgfpointdiff{\pgfpointanchor{SNsidebox}{south}}{\pgfpointanchor{SNsidebox}{north}}}\typeout{SNHEIGHT:7:9:\the\SNsideheight}}
\expandafter\gdef\csname SNside8\endcsname{\node[anchor=north west,inner sep=0pt,text width=48mm,font=\SideFont,align=left] (SNsidebox) at ([xshift=138mm,yshift=-95.00pt]current page.north west) {\textcolor{SNink}{\textbf{术语 · 网络的连接对象}}\par\vspace{1.5mm}网络连接设备，互连网连接多个网络；Internet采用TCP/IP协议族。};\pgfextracty{\SNsideheight}{\pgfpointdiff{\pgfpointanchor{SNsidebox}{south}}{\pgfpointanchor{SNsidebox}{north}}}\typeout{SNHEIGHT:8:11:\the\SNsideheight}}
\expandafter\gdef\csname SNside9\endcsname{\node[anchor=north west,inner sep=0pt,text width=48mm,font=\SideFont,align=left] (SNsidebox) at ([xshift=138mm,yshift=-95.00pt]current page.north west) {\textcolor{SNink}{\textbf{不变量 · 出口何时开始}}\par\vspace{1.5mm}下一跳既要等本组收齐并处理，也要等出口链路空闲，取两个时刻的较大值。};\pgfextracty{\SNsideheight}{\pgfpointdiff{\pgfpointanchor{SNsidebox}{south}}{\pgfpointanchor{SNsidebox}{north}}}\typeout{SNHEIGHT:9:13:\the\SNsideheight}}
\expandafter\gdef\csname SNside10\endcsname{\node[anchor=north west,inner sep=0pt,text width=48mm,font=\SideFont,align=left] (SNsidebox) at ([xshift=138mm,yshift=-95.00pt]current page.north west) {\textcolor{SNink}{\textbf{不变量 · 回忆时说出计时起点}}\par\vspace{1.5mm}同一网络从请求建立连接、首次发送数据或收到第一位开始计时，会得到不同表达式。主动回忆时先口述起点与终点，再写公式；把条件也回忆出来，才可能迁移到变式题。};\pgfextracty{\SNsideheight}{\pgfpointdiff{\pgfpointanchor{SNsidebox}{south}}{\pgfpointanchor{SNsidebox}{north}}}\typeout{SNHEIGHT:10:10:\the\SNsideheight}}
\expandafter\gdef\csname SNside11\endcsname{\node[anchor=north west,inner sep=0pt,text width=48mm,font=\SideFont,align=left] (SNsidebox) at ([xshift=138mm,yshift=-95.00pt]current page.north west) {\textcolor{SNink}{\textbf{课程连接 · 原语与报文}}\par\vspace{1.5mm}服务原语是同机层间事件；一次原语可引起多个对等协议报文。};\pgfextracty{\SNsideheight}{\pgfpointdiff{\pgfpointanchor{SNsidebox}{south}}{\pgfpointanchor{SNsidebox}{north}}}\typeout{SNHEIGHT:11:14:\the\SNsideheight}}
\AddToShipoutPictureFG{\begin{tikzpicture}[remember picture,overlay]\ifcsname SNside\thepage\endcsname\csname SNside\thepage\endcsname\fi\end{tikzpicture}}